\documentclass[11pt]{article}
\usepackage[margin=1in]{geometry}
\usepackage{amsmath,amssymb}
\usepackage[hidelinks]{hyperref}

\title{Convergent denominators of \(e\)\\[2pt]\large Proof of Conjecture 00000000346}
\author{}
\date{}

\begin{document}
\maketitle

\noindent
Let $q_n$ be the denominator of the $n$th regular continued-fraction
convergent of $e$, with the integer part indexed by zero. For every real
$\varepsilon>0$ we prove the limit
\[
\lim_{n\to\infty}\frac{q_{n+1}}{q_n^{1+\varepsilon}}=0,
\]
which gives the asserted limit superior.

\section*{Denominator bounds}

Euler's continued fraction for $e$~\cite[\S21]{euler} has $a_0=2$ and,
for every integer $j\ge 0$,
\[
(a_{3j+1},a_{3j+2},a_{3j+3})=(1,2(j+1),1).
\]
Cohn~\cite[Theorem~1]{cohn} gives a modern proof in the equivalent form
that starts $[1,0,1,\ldots]$, where the initial $1,0,1$ combines to $2$.
The convergent denominators obey $q_{-1}=0$, $q_0=1$, and
\[
q_{n+1}=a_{n+1}q_n+q_{n-1}\qquad(n\ge0).
\]
The block formula gives $1\le a_{n+1}\le n+1$. The recurrence makes
the denominators positive and nondecreasing for $n\ge0$. Consequently,
\[
q_{n+1}\le(n+2)q_n,
\qquad
q_{n+2}\ge q_{n+1}+q_n\ge2q_n,
\qquad
q_{2s}\ge2^s\quad(s\ge0),
\]
where the last inequality follows by induction from $q_0=1$.

\section*{A bound using integer powers}

Fix positive integers $r,B$, set $d=4r$, and put $s=\lfloor n/d\rfloor$.
For $n\ge B(d+2)d$, we have $s\ge B(d+2)$. As $n<d(s+1)$,
\[
B(n+2)\le B(d+2)(s+1)\le s(s+1)<(s+1)^2.
\]
The inequality $s+1\le2^s$ follows by induction on $s$. Since $n\ge ds$,
the denominator bound gives
\[
q_n\ge q_{ds}\ge 2^{2rs}=\bigl((2^s)^2\bigr)^r
    \ge \bigl((s+1)^2\bigr)^r.
\]
It follows that $(B(n+2))^r<q_n$, and hence
\begin{equation}
 (Bq_{n+1})^r\le (B(n+2))^r q_n^r<q_n^{r+1}.
 \label{eq:integer-bound}
\end{equation}

\section*{Passage to real exponents}

Let $\varepsilon>0$ and $\delta>0$. By the Archimedean property, choose
positive integers $r$ and $B$ with $1/r\le\varepsilon$ and $1/B<\delta$.
Taking the positive $r$th root of~\eqref{eq:integer-bound} and dividing by
$Bq_n^{1+1/r}$ gives, for all $n\ge B(d+2)d$,
\[
0\le \frac{q_{n+1}}{q_n^{1+\varepsilon}}
\le \frac{q_{n+1}}{q_n^{1+1/r}} < \frac1B<\delta,
\]
using $q_n\ge1$ for the middle inequality. This proves the limit.

The Lean 4 project formalizes the displayed block formula and the
denominator recurrence. Its theorem
\texttt{C346.eventual\_rational\_bounds} proves~\eqref{eq:integer-bound}
for every positive $r$ and $B$, using the threshold above. Run
\texttt{lake build} in the \texttt{lean} directory to check the integer proof.

\begin{thebibliography}{2}
\bibitem{euler}
L.~Euler, \emph{De fractionibus continuis dissertatio}, Commentarii
Academiae Scientiarum Petropolitanae \textbf{9} (1744), 98--137, E71.
\href{https://scholarlycommons.pacific.edu/euler-works/71/}{Euler Archive}.

\bibitem{cohn}
H.~Cohn, \emph{A short proof of the simple continued fraction expansion of $e$},
American Mathematical Monthly \textbf{113} (2006), 57--62.
\href{https://doi.org/10.2307/27641837}{DOI 10.2307/27641837}.
\end{thebibliography}

\end{document}
